\section{Introduction}
\label{sec:intro}
Autoregressive modeling supports incremental construction and completion of three-dimensional scenes. GaussianGPT~\cite{gaussiangpt} compresses 3D Gaussians into a sparse discrete representation and alternates spatial occupancy positions with their associated scene features. Thousands of output tokens make serial Transformer calls a major latency source.

Speculative decoding amortizes target execution by verifying cheaply proposed tokens in parallel~\cite{leviathan}. Recent mesh-generation methods also exploit representation structure: FlashMesh uses mesh hierarchy and vertex sharing, while XSpecMesh combines multiple draft heads with target verification~\cite{flashmesh,xspecmesh}. For GaussianGPT, the design question is how to assign draft computation across two token types with different cheaply available conditioning information.

The P/F representation supplies two complementary resources. Legal spatial order and position history support a small ranked set of position proposals. Feature prediction can reuse an available target state and condition on a proposed position. This permits both positions in a short P/F path to be selected before feature drafting, while a single neural prediction directly estimates the state needed for the intervening feature. The target verifies every token under its complete history. Proposal-side simplification therefore changes the computation required to offer a path while retaining exact target sampling through acceptance and residual resampling.

We ask: \emph{How should these resources be organized to increase useful generation advance per unit of execution cost?} \PF\ couples a Position Ranker with a conditional Feature MLP. The ranker selects legal positions, and the MLP predicts a candidate-position state residual from an existing target state, the previous feature, and the candidate position. A shared target head supplies the feature proposal. Their jointly verified P/F path avoids intermediate neural draft-state updates and a separate cross-round draft history cache. This allocation couples spatial ranking with target-state-assisted prediction.

Our experiments trace this allocation from position-only speculation to the feature increment, state-input ablations, and wider or heavier drafts. Wider or heavier drafts can increase advance or acceptance while reducing measured throughput. Analysis of the same path exposes a remaining serial dependency: after position rejection, the target must process the residual-sampled position to generate its feature. Position Rescue verifies one sibling position in advance and reuses its conditional state when selected by the original residual sampler. It complements the main path by reducing recovery work without changing reference acceptance events.

\PF\ achieves approximately $1.58\times$ throughput over a target with shared execution optimizations. The basic P/F path also accelerates complete, natural-EOS single-chunk requests over thousands of tokens; reference-image statistics and rendered outputs accompany these results. The current full-method natural-EOS increment remains unmeasured. Our contributions are:
\begin{itemize}[leftmargin=*,itemsep=2pt]
\item \textbf{A P/F-specific allocation of draft computation.} Legal position ranking and direct candidate-state prediction form a lightweight, jointly verifiable short path. We evaluate how feature prediction extends the measured position-only gain and how state information and draft cost affect that increment.
\item \textbf{Recovery-state reuse within the same P/F decoding process.} A single backup position avoids serial recovery on a residual hit. We distinguish useful advance from avoided recovery work, state the preservation conditions, and evaluate both decoding behavior and complete single-chunk outputs.
\end{itemize}
